% authority: science_draft
\documentclass[11pt]{article}
\usepackage[margin=1in]{geometry}
\usepackage{amsmath,amssymb,amsthm}
\usepackage{longtable}

\newcommand{\GSC}{\mathsf{GSC}}
\newcommand{\Msrc}{M_{\mathrm{src}}}
\newcommand{\Mgsccand}{M_{\mathrm{src}}^{\GSC\mbox{-}\mathrm{cand}}}
\newcommand{\Qsrc}{Q_{\mathrm{src}}^{\GSC}}
\newcommand{\Bottomsrc}{\mathrm{Bottom}_{\mathrm{src}}}
\newcommand{\Covsrc}{\mathrm{Cov}_{\mathrm{src}}^{\GSC}}
\newcommand{\geff}{g_{\mathrm{eff}}}

\newtheorem{definition}{Definition}
\newtheorem{lemma}{Lemma}
\newtheorem{condition}{Condition}
\newtheorem{rulebox}{Rule}

\title{Source Transition-Token Groupoid Conditions for \(M_{\mathrm{src}}^{\mathsf{GSC}\mbox{-}\mathrm{cand}}(E)\)}
\author{Ontology Formalizer Packet RT-20260614-126}
\date{2026-06-24}

\begin{document}
\maketitle

\section{Control Status}

This artifact is the Phase 3 Ontology Formalizer output for
\texttt{RT-20260614-126}. The selected role is
\texttt{ontology-formalizer@0.2.0}. The parent handoff is
\texttt{handoff-0161}. The target derivation milestone is
\texttt{source\_manifold\_m\_src}.

This packet formalizes transition-token candidate conditions over the selected
Phase 2 site-coverage route. It does not adopt a transition groupoid, does not
adopt full \(\Msrc\), does not construct full \(\Msrc\), does not edit
canonical ontology TeX, does not define \(\geff\), does not derive matter
coupling, does not derive Einstein equations, does not promote benchmark
status, does not issue Gate Chair closure, does not claim completed
derivation, does not claim future source-extension impossibility, and does
not reject the global theory.

The input candidate remains \(\Mgsccand(E)\) as reviewed source-extension
bridge-candidate data. The selected topology analogue remains
\(\Covsrc(E)\), a source-only site or coverage route. The result of this
packet is:
\[
  \texttt{transition\_groupoid\_status = formalized\_candidate\_conditions}.
\]

\section{Inputs}

The candidate data are
\[
  \Mgsccand(E)=
  \left(
  \Qsrc(E),
  I_E,
  \{U_i^{\mathrm{src}}\},
  \{O_{ij}^{\mathrm{src}}\},
  \{\chi_i^{\mathrm{src}}\},
  \{\tau_{ij}^{\mathrm{src}}\},
  \mathrm{Cert}(E),
  \Bottomsrc
  \right).
\]

Phase 2 selected \(\Covsrc(E)\) as the source-only topology analogue route.
Here \(\Covsrc(E)\) is used only as a coverage domain on which source overlap
and transition-token obligations can be stated. It is not a target topology,
not a target smooth atlas, and not adopted \(\Msrc\).

\section{Source Overlap Domains}

\begin{definition}[Source overlap domain]
For \(i,j\in I_E\), define the source overlap domain as the source predicate
\[
  D_{ij}^{\mathrm{src}}(E)
  :=
  O_{ij}^{\mathrm{src}}
  \subseteq_{\mathrm{src}}
  U_i^{\mathrm{src}}\times U_j^{\mathrm{src}},
\]
when the selected coverage route allows subset notation. If subset notation
is not source-defined, \(D_{ij}^{\mathrm{src}}(E)\) denotes the binary source
predicate
\[
  D_{ij}^{\mathrm{src}}(E)(u_i,u_j)
  \equiv
  O_{ij}^{\mathrm{src}}(u_i,u_j).
\]
\end{definition}

The triple-overlap support is the source predicate
\[
  D_{ijk}^{\mathrm{src}}(E)
  :=
  D_{ij}^{\mathrm{src}}(E)
  \wedge D_{jk}^{\mathrm{src}}(E)
  \wedge D_{ik}^{\mathrm{src}}(E),
\]
read only as coverage-compatible source support. It is not a target triple
intersection of coordinate charts.

\section{Transition Tokens}

\begin{definition}[Source transition token]
For each source overlap domain with non-bottom support, a transition token is
a symbolic source arrow
\[
  \tau_{ij}^{\mathrm{src}}:
  D_{ij}^{\mathrm{src}}(E)\rightsquigarrow D_{ji}^{\mathrm{src}}(E).
\]
The symbol \(\rightsquigarrow\) records source-transition syntax only. It
does not assert a function space, differentiability class, target coordinate
transition map, or smooth atlas.
\end{definition}

\begin{definition}[Identity token]
An identity token \(id_i^{\mathrm{src}}\) is source-defined when
\(\mathrm{Cert}(E)\) or the selected coverage route supplies a non-bottom
identity certificate on \(D_{ii}^{\mathrm{src}}(E)\). When present, write
\[
  \tau_{ii}^{\mathrm{src}} = id_i^{\mathrm{src}}.
\]
If no source identity token is supplied, the branch returns \(\Bottomsrc\) or
the primary obstruction \texttt{OB-MSRC-TRANSITION-INVERSE-FAIL}.
\end{definition}

\begin{definition}[Source composition predicate]
The partial source composition predicate
\[
  \mathrm{Comp}_{\mathrm{src}}
  \left(\tau_{jk}^{\mathrm{src}},\tau_{ij}^{\mathrm{src}}\right)
\]
is defined only on coverage-compatible triple support
\(D_{ijk}^{\mathrm{src}}(E)\) and only when source certificates supply a
non-bottom composition witness. If defined, the composite is written
\[
  \tau_{jk}^{\mathrm{src}}\circ_{\mathrm{src}}\tau_{ij}^{\mathrm{src}}.
\]
The composition symbol is source syntax. It does not import target coordinate
composition.
\end{definition}

\section{Lemma Schemas}

\begin{lemma}[Identity candidate condition]
Assume \(D_{ii}^{\mathrm{src}}(E)\) is non-bottom and
\(\mathrm{Cert}(E)\) supplies an identity certificate for support \(i\). Then
the candidate transition-token system satisfies the identity schema
\[
  \tau_{ii}^{\mathrm{src}} = id_i^{\mathrm{src}}.
\]
If the identity certificate is absent or degenerate, the branch returns
\(\Bottomsrc\) or \texttt{OB-MSRC-TRANSITION-INVERSE-FAIL}.
\end{lemma}

\begin{lemma}[Inverse candidate condition]
Assume \(D_{ij}^{\mathrm{src}}(E)\) and \(D_{ji}^{\mathrm{src}}(E)\) are
non-bottom and \(\mathrm{Cert}(E)\) supplies source inverse certificates for
\(\tau_{ij}^{\mathrm{src}}\) and \(\tau_{ji}^{\mathrm{src}}\). Then the
candidate transition-token system satisfies
\[
  \tau_{ji}^{\mathrm{src}}\circ_{\mathrm{src}}\tau_{ij}^{\mathrm{src}}
  =
  id_i^{\mathrm{src}},
  \qquad
  \tau_{ij}^{\mathrm{src}}\circ_{\mathrm{src}}\tau_{ji}^{\mathrm{src}}
  =
  id_j^{\mathrm{src}},
\]
on the corresponding source overlap supports. If inverse support or source
composition is absent or degenerate, the branch returns \(\Bottomsrc\) or
\texttt{OB-MSRC-TRANSITION-INVERSE-FAIL}.
\end{lemma}

\begin{lemma}[Cocycle candidate condition]
Assume \(D_{ijk}^{\mathrm{src}}(E)\) is non-bottom and
\(\mathrm{Cert}(E)\) supplies source composition and cocycle certificates.
Then the candidate transition-token system satisfies one of the equivalent
source cocycle schemas chosen by the later integrated theorem:
\[
  \tau_{ki}^{\mathrm{src}}\circ_{\mathrm{src}}
  \tau_{jk}^{\mathrm{src}}\circ_{\mathrm{src}}
  \tau_{ij}^{\mathrm{src}}
  =
  id_i^{\mathrm{src}},
\]
or
\[
  \tau_{jk}^{\mathrm{src}}\circ_{\mathrm{src}}\tau_{ij}^{\mathrm{src}}
  =
  \tau_{ik}^{\mathrm{src}}.
\]
If triple support, source composition, or source cocycle certificates are
absent or degenerate, the branch returns \(\Bottomsrc\) or
\texttt{OB-MSRC-TRANSITION-COCYCLE-FAIL}.
\end{lemma}

\section{Bottom and Obstruction Rule}

\begin{rulebox}[Transition fail-closed rule]
If identity, inverse, composition, or cocycle structure cannot be stated from
source-local supports, source overlap predicates, source transition tokens,
\(\mathrm{Cert}(E)\), and \(\Bottomsrc\) guards, the branch must return
\(\Bottomsrc\) or exactly one primary obstruction label. The allowed primary
labels for this packet are:
\[
\begin{array}{ll}
\texttt{OB-MSRC-TRANSITION-INVERSE-FAIL}, &
\text{identity or inverse support absent, nonunique, or degenerate},\\
\texttt{OB-MSRC-TRANSITION-COCYCLE-FAIL}, &
\text{composition or triple-overlap coherence absent or degenerate},\\
\texttt{OB-MSRC-BOTTOM-LEAK}, &
\text{failure branch becomes non-bottom by notation or process authority},\\
\texttt{OB-MSRC-ADOPTION-THEOREM-UNFORMED}, &
\text{transition conditions cannot be stated without a new primitive}.
\end{array}
\]
\end{rulebox}

\section{No-Target Guard}

The transition tokens \(\tau_{ij}^{\mathrm{src}}\) are not target
coordinate-transition maps and do not inherit smoothness from a target atlas.
The following arguments are forbidden:

\begin{itemize}
  \item ``They look like chart transitions.''
  \item ``A manifold atlas would require this, so assume it.''
  \item ``The exact-GR benchmark supplies smooth transitions.''
  \item ``The target atlas is the intended model.''
  \item ``Validator PASS confirms transition coherence.''
\end{itemize}

Allowed evidence is narrower:

\begin{itemize}
  \item \(\mathrm{Cert}(E)\) contains source identity, inverse, composition,
  or cocycle certificates.
  \item The selected \(\Covsrc(E)\) route provides a source-only support
  domain on which composition can be stated.
  \item Failure to compose returns \(\Bottomsrc\) or a named obstruction.
  \item Future finite checker counterexamples may demonstrate where a
  condition fails, but checker PASS is not adoption proof.
\end{itemize}

\section{Adoption Theorem Effect}

This packet narrows assumption A4,
\(\mathrm{transition\_inverse\_cocycle}(E)\), from a vague obligation to a
source transition-token candidate-condition family. It does not discharge A4
as an adopted fact. It also does not discharge carrier support, chart-support
regularity, separation, countability, soldering, finite-variation robustness,
bottom behavior, no-target-import, or the local-to-global representability
burden.

\begin{longtable}{|p{0.26\linewidth}|p{0.25\linewidth}|p{0.39\linewidth}|}
\hline
\textbf{Obligation} & \textbf{Status after Phase 3} & \textbf{Reason} \\
\hline
A2 topology analogue & selected route & \(\Covsrc(E)\) is selected from Phase 2, formal proof remains pending. \\
\hline
A4 transition inverse & formalized candidate condition & Identity and inverse schemas are stated with source certificates and bottom branches. \\
\hline
A4 transition cocycle & formalized candidate condition & Composition and triple-overlap schemas are stated with source certificates and bottom branches. \\
\hline
A8 bottom behavior & narrowed & Missing or degenerate transition data returns \(\Bottomsrc\) or an obstruction. \\
\hline
A9 no-target-import & narrowed but unaudited & Target coordinate maps and smooth atlas imports are explicitly forbidden. Later audit remains required. \\
\hline
\end{longtable}

\section{Next Route}

The next legal packet is Phase 4: one bounded finite/local witness
classification packet. That packet should separate finite or local transition
and coverage survival from global \(\Msrc\) adoption. It must not route to
\(\geff\), matter coupling, Einstein equations, benchmark promotion, Gate
Chair closure, completed derivation, future source-extension impossibility,
or global theory rejection.

\section{Conclusion}

The transition-token groupoid candidate conditions for \(\Mgsccand(E)\) are
now formalized over the selected source-only site-coverage route. The
adoption-level obligation remains open. Same-milestone continuation remains
open. Downstream GR claims remain blocked. No target coordinate-transition
map is used as mathematical evidence.

\begin{thebibliography}{9}

\bibitem{plan}
The AEther-Flow Research Project. (2026, June 24). \emph{Recommendations
implementation plan for continue-research v4} [Internal implementation plan].

\bibitem{theorem-envelope}
The AEther-Flow Research Project. (2026, June 24). \emph{Minimum source-only
adoption theorem envelope for M\_src GSC candidate} [Internal
research-control artifact].

\bibitem{route-decision}
The AEther-Flow Research Project. (2026, June 24). \emph{Source-only
topology analogue route decision for M\_src GSC candidate} [Internal
research-control artifact].

\bibitem{bridge-attempt}
The AEther-Flow Research Project. (2026, June 23). \emph{GenCover-to-Msrc
bridge attempt} [Internal research-control artifact].

\bibitem{gate-review}
The AEther-Flow Research Project. (2026, June 23). \emph{Resp\_lc M\_src GSC
bridge candidate human-gate adoption review} [Internal research-control
artifact].

\end{thebibliography}

\end{document}
